\section{Introduction}
\label{sec:introduction}

\subsection{Economic motivation and forecasting problem}
\label{subsec:intro_motivation}

Assessing the current state of the economy is inherently a real-time forecasting
problem. Policymakers, financial-market participants, firms, and researchers
frequently need estimates of output growth, inflation, and labour-market
conditions before the corresponding official statistics are fully observed.
Macroeconomic indicators are published at different frequencies and on different
release schedules, so the information available near the beginning of a target
period can differ substantially from that available only a few weeks later.
Real-time macroeconomic datasets consequently have a ragged or ``jagged'' edge,
and nowcasts are naturally revised as new observations become available
\citep{giannone2008nowcasting,banbura2011blackbox}. A real-time forecast should
therefore be regarded as conditional not only on a target period but also on the
stage of the information flow. Denoting by $\It_{t,s}$ the information available
for target period $t$ at release stage $s$, two stages within the same target
period generally satisfy
\begin{equation}
    \It_{t,s_1} \neq \It_{t,s_2},
    \qquad s_1 \neq s_2.
    \label{eq:stage_information_sets}
\end{equation}
Here, $s$ represents a position in the release calendar rather than a conventional
forecast horizon. This distinction is central to the empirical design because
forecasts generated at different stages need not be based on the same observed
variables, even when they refer to the same macroeconomic outcome.

The difficulty is compounded by instability in the economic relationships used
for prediction. Macroeconomic forecasting equations have long been found to
display substantial instability, while predictive performance can vary even
when conventional tests do not identify a single discrete structural break
\citep{stockwatson1996instability,rossi2021instability}. The targets considered
in this paper provide a particularly demanding setting in this respect.
Descriptive evidence from the study sample illustrates the scale of the
variation around the COVID-19 episode. The standard deviation of annualised
monthly core CPI inflation rises from approximately 0.74 percentage points over
2015--2019 to about 2.50 percentage points thereafter, while the standard
deviation of monthly payroll changes increases from roughly 83 thousand workers
to about 2.44 million, the latter being heavily influenced by the extraordinary
pandemic labour-market observations. Annualised quarter-on-quarter real GDP
growth, meanwhile, fell to approximately $-32.8\%$ in 2020Q2 before rebounding
to approximately $29.9\%$ in 2020Q3. These observations are descriptive rather
than evidence that any particular forecasting method should dominate. They
instead illustrate why assuming a single forecasting relationship to perform
uniformly across targets, periods, and economic environments is a demanding
restriction.

Variation in the information set creates a second and conceptually distinct
source of possible forecast instability. Incoming releases do not merely add
observations mechanically; their forecasting value depends on when they become
available and on what has already been observed. Accounting explicitly for
publication lags can materially change the measured usefulness of individual
predictors, with relatively timely survey or financial indicators potentially
contributing more at early forecast origins before harder activity data are
released \citep{banbura2011blackbox}. Moreover, real-time evidence for U.S. GDP
shows that the relative ranking of forecasting models can change as information
accumulates during the quarter \citep{aastveit2014density}. This does not imply
that changing models across stages must improve forecasting accuracy. Indeed,
whether such variation can be exploited successfully out of sample is itself an
empirical question. It does, however, imply that a comparison between a fixed
forecasting rule and a rule conditioned on the stage of the release calendar is
economically meaningful. Importantly, such a comparison must distinguish the
effect of receiving additional information from the separate effect of changing
the forecasting model: otherwise an improvement at a later stage could simply
reflect the richer information set rather than any gain from stage-dependent
model choice.

Real-time evaluation also requires distinguishing the release-calendar problem
from the data-vintage problem. Macroeconomic observations available in current
databases may differ from the values that were actually available to a
forecaster at a historical forecast origin. Revisions can therefore alter both
the dependent variable and the predictor histories used in retrospective
forecast exercises, which is the central motivation for constructing
vintage-specific real-time datasets \citep{croushorestark2001realtime,
croushore2011frontiers}. In the notation used here, the historically feasible
information set $\It_{t,s}^{RT}$ need not coincide with a reconstruction based
on the latest revised observations, $\It_{t,s}^{LV}$. The relevance of this
distinction is visible even in a simple validation exercise using selected
ALFRED snapshots. Over 2010Q1--2018Q3, the absolute difference between
annualised quarterly GDP growth calculated from the January 2019 and January
2020 vintages averages approximately 0.22 percentage points and reaches about
1.22 percentage points; for 2017Q4 alone, the corresponding growth estimate
changes from roughly $2.27\%$ to $3.48\%$. These figures are only illustrative
differences between selected annual vintages: they should not be interpreted as
an estimated revision distribution or as a first-release-versus-final-release
experiment. Their role is instead to demonstrate why a historically feasible
forecast experiment and a latest-vintage counterfactual are not necessarily the
same exercise.

Taken together, sequential information arrival, changing data vintages, and
instability in predictive relationships transform the forecasting problem from
the selection of one universally preferred model into a comparison of
forecasting rules conditional on the information actually available at each
forecast origin. Model uncertainty also makes forecast combinations an
important benchmark, since pooling can provide robustness when the identity of
the best individual specification varies over time
\citep{timmermann2006combinations}. Existing real-time GDP research already
shows both within-quarter variation in model rankings and strong performance
from forecast combinations \citep{aastveit2014density}; consequently, the
empirical question in this paper is not whether stage dependence or model
adaptation is intrinsically novel or necessarily superior. Rather, the question
is whether, within a common pseudo-real-time design spanning GDP, inflation,
and labour-market targets, the predictive value of information changes
systematically across release stages, whether a pre-specified
stage-dependent forecasting policy adds value beyond that information-arrival
effect, and whether the resulting conclusions are sensitive to the use of
historically available rather than latest-revised data. These questions motivate
the formal research questions and hypotheses developed in the following
subsections.

\subsection{Research questions}
\label{subsec:research_questions}

The central research question is whether macroeconomic nowcasting performance can be improved by explicitly conditioning model choice and predictive uncertainty on the stage of the real-time information set.

The empirical analysis addresses four linked questions:
\begin{enumerate}[label=\textbf{RQ\arabic*:}, leftmargin=*]
    \item Does incremental information arriving within the target period improve nowcast accuracy when the forecasting model is held fixed?
    \item Does a pre-specified stage-dependent model policy add predictive value beyond the information-arrival effect captured in RQ1?
    \item Do historically available data vintages materially alter measured forecast performance relative to latest-vintage data?
    \item Do prediction intervals calibrated only from previously observable forecast errors achieve appropriate coverage while remaining sufficiently sharp?
\end{enumerate}

\subsection{Research hypotheses and empirical decision rules}
\label{subsec:hypotheses}

The four research questions are translated into the following pre-specified hypotheses. Throughout the paper, inference is stated in terms of rejection or failure to reject the corresponding null hypothesis rather than ``acceptance'' of an alternative.

\paragraph{H1: Incremental real-time information has predictive value}
Let $e_{k,t,s}^{(m)}$ denote the forecast error for target $k$ from a \emph{fixed} model $m$ at information stage $s$. For adjacent stages $s$ and $s+1$, define the loss differential
\begin{equation}
    d_{k,t}^{H1}(m;s,s+1)
    =
    L\!\left(e_{k,t,s+1}^{(m)}\right)
    -
    L\!\left(e_{k,t,s}^{(m)}\right).
    \label{eq:h1_loss_diff}
\end{equation}
The null and alternative are
\begin{align}
    H_{0}^{(1)} &: \E\!\left[d_{k,t}^{H1}\right] \geq 0,
    \label{eq:h1_null}\\
    H_{A}^{(1)} &: \E\!\left[d_{k,t}^{H1}\right] < 0.
    \label{eq:h1_alt}
\end{align}
The model is deliberately held fixed across stages so that H1 isolates the marginal predictive value of the expanding real-time information set. The primary inferential loss is squared error, with absolute error used as a secondary loss metric; RMSE and MAE provide descriptive summaries in the corresponding units. Statistical evidence is based on paired loss differentials with two-sided Newey--West heteroskedasticity-and-autocorrelation-consistent inference.

\paragraph{H2: Stage-dependent model policies outperform fixed-model forecasting}
Let $\mathcal{P}(k,s)=m^{*}_{k,s}$ denote the pre-specified stage-dependent policy, and let $m_k^{F}$ denote a fixed model selected for target $k$ using a policy-development sample only. Define
\begin{equation}
    d_{k,t}^{H2}
    =
    L\!\left(e_{k,t}^{\mathrm{stage}}\right)
    -
    L\!\left(e_{k,t}^{\mathrm{fixed}}\right).
    \label{eq:h2_loss_diff}
\end{equation}
The null and alternative are
\begin{align}
    H_{0}^{(2)} &: \E\!\left[d_{k,t}^{H2}\right] \geq 0,
    \label{eq:h2_null}\\
    H_{A}^{(2)} &: \E\!\left[d_{k,t}^{H2}\right] < 0.
    \label{eq:h2_alt}
\end{align}
H2 is the principal hypothesis of the paper. It tests whether allowing model choice to depend on the forecast stage provides additional predictive value after the changing information set itself has been accounted for. To avoid policy-selection bias, the stage policy and the fixed comparator are selected using observations prior to the final evaluation sample. Inference uses paired squared-error differentials, with absolute-error differentials as a secondary metric and RMSE/MAE as descriptive summaries.

For H1 and H2, the economic hypotheses are directional, but the archived
confirmatory implementation reports conservative two-sided HAC \(p\)-values;
directional support is therefore determined jointly by the sign of the loss
differential and the two-sided test.

\paragraph{H3: Real-time data vintages materially affect forecast evaluation}
For identical forecast origins, models, and stage definitions, let $e_{k,t}^{RT}$ denote forecast errors generated from historically available vintages and $e_{k,t}^{LV}$ denote errors from a latest-vintage counterfactual. Define
\begin{equation}
    d_{k,t}^{H3}
    =
    L\!\left(e_{k,t}^{LV}\right)
    -
    L\!\left(e_{k,t}^{RT}\right).
    \label{eq:h3_loss_diff}
\end{equation}
The hypothesis is intentionally two-sided:
\begin{align}
    H_{0}^{(3)} &: \E\!\left[d_{k,t}^{H3}\right] = 0,
    \label{eq:h3_null}\\
    H_{A}^{(3)} &: \E\!\left[d_{k,t}^{H3}\right] \neq 0.
    \label{eq:h3_alt}
\end{align}
The confirmatory analysis examines changes in forecast loss and forecast values for the frozen stage policy. With the loss differential defined as latest-vintage minus real-time loss, negative values indicate lower loss under the latest-vintage counterfactual. Candidate-model ranking sensitivity is outside the present analysis. The sign of the vintage effect is left to the data rather than imposed ex ante.

\paragraph{H4: Prior-only empirical intervals are calibrated and sharp}
For an 80\% predictive interval $[L_t,U_t]$, define the violation indicator
\begin{equation}
    V_t = \I\!\left(y_t<L_t \;\text{or}\; y_t>U_t\right).
    \label{eq:h4_violation}
\end{equation}
The calibration component of H4 is assessed against nominal coverage:
\begin{align}
    H_{0}^{(4a)} &: \Prob(L_t\leq y_t\leq U_t)=0.80,
    \label{eq:h4_coverage_null}\\
    H_{A}^{(4a)} &: \Prob(L_t\leq y_t\leq U_t)\neq0.80.
    \label{eq:h4_coverage_alt}
\end{align}
In addition, interval violations should not exhibit systematic serial dependence. Because correct coverage can be achieved with uninformatively wide intervals, calibration is evaluated jointly with sharpness using average interval width and the proper interval score. The prior-only calibration method is compared with conventional and alternative empirical interval benchmarks.

\begin{table}[H]
\centering
\caption{Pre-specified research hypotheses and primary empirical evidence}
\label{tab:hypotheses}
\begin{threeparttable}
\footnotesize
\setlength{\tabcolsep}{3pt}
\begin{tabular}{>{\raggedright\arraybackslash}p{0.055\textwidth}>{\raggedright\arraybackslash}p{0.285\textwidth}>{\raggedright\arraybackslash}p{0.285\textwidth}>{\raggedright\arraybackslash}p{0.265\textwidth}}
\toprule
Hyp. & Scientific question & Null hypothesis & Primary evidence \\
\midrule
H1 & Does incremental information improve nowcasts when the model is fixed? & Later-stage expected loss is no lower than earlier-stage expected loss. & RMSE/MAE changes; paired loss-differential test. \\
H2 & Does stage-dependent model choice add value beyond information arrival? & Expected loss of the stage policy is no lower than that of a fixed policy. & Relative RMSE; MAE; predictive-accuracy tests. \\
H3 & Do historically available vintages change forecast evaluation? & Real-time and latest-vintage expected losses are equal. & Loss differential; forecast-value revisions. \\
H4 & Are prior-only 80\% intervals calibrated and informative? & Coverage equals the nominal rate and violations show no systematic dependence. & Coverage; independence; width; interval score. \\
\bottomrule
\end{tabular}
\begin{tablenotes}[flushleft]\footnotesize
\item Notes: H1 holds the forecasting model fixed across information stages. H2 then isolates the incremental value of stage-dependent model choice. H3 is a two-sided validity test. H4 requires assessment of both calibration and sharpness rather than coverage alone.
\end{tablenotes}
\end{threeparttable}
\end{table}

\subsection{Empirical scope}
\label{subsec:empirical_scope}

The framework is evaluated across three macroeconomic domains: current-quarter real GDP growth; monthly headline CPI, core CPI, headline PCE, and core PCE inflation; and monthly payroll employment, unemployment, and average hourly earnings. These targets differ in frequency, persistence, revision behaviour, publication timing, and the amount of within-period information available, providing a heterogeneous test bed for the four hypotheses.

\subsection{Contributions}
\label{subsec:contributions}

The paper makes four contributions. First, it reconstructs historically feasible information sets at each forecast origin rather than evaluating models on a fully revised rectangular dataset. Second, it formalises stage-dependent forecast policies that allow the preferred forecasting rule to vary with the information available within the target period. Third, it evaluates the same general architecture across GDP, inflation, and labour-market targets with distinct release calendars and information structures. Fourth, it evaluates predictive intervals calibrated only from forecast errors that were observable before the forecast origin.

\subsection{Main findings}
\label{subsec:main_findings}

The confirmatory evidence is heterogeneous across hypotheses and domains.
For H1, later information stages lower mean squared loss in 16 of 28
adjacent-stage comparisons, but only four improvements and two deteriorations
are significant at the 5\% level, so information arrival is valuable without
being monotonically beneficial.  For H2, the frozen stage policy differs from
the fixed comparator in 11 target--stage cells; seven have lower mean squared
loss, and four statistically significant gains occur in labour-market cells
(AHE at month end and pre-report, PAYEMS at month open, and UNRATE at month
open), with no significant stage-policy loss.  H3 shows that data-vintage
effects are likewise target specific: latest-vintage data substantially
improve unemployment evaluation, whereas core PCE absolute-error evaluation
moves significantly in the opposite direction and most other target-level
effects are statistically weaker.  Finally, H4 finds that prior-only empirical
interval calibration often improves interval scores for labour and core CPI,
while GDP's 100\% empirical coverage reflects conservative overcoverage rather
than ideal 80\% calibration.  Taken together, the results favour
stage-aware and vintage-disciplined forecast evaluation, but not a single
forecasting or interval rule that dominates uniformly across targets and
release stages. The predeclared robustness analysis preserves this qualitative
interpretation: shortening H2 policy development changes only 2 of 36 stage
mappings, the stage-level H3 decomposition shows that the UNRATE vintage effect
is broad across the release cycle, and Holm correction narrows the set of
individually significant point-forecast findings. H4 is more tuning-sensitive
for GDP, where alternative rolling calibration windows improve interval scores
relative to the frozen 20-observation rule.

\subsection{Structure of the paper}
\label{subsec:paper_structure}

The remainder of the paper is organised as follows. \Cref{sec:literature} reviews the related literature. \Cref{sec:data} defines the data and real-time information sets. \Cref{sec:methodology} presents the forecasting models and stage-dependent policy framework. \Cref{sec:experiment} describes the pseudo-real-time forecast experiment and the tests associated with H1--H4. \Cref{sec:results} reports the empirical results in the same hypothesis order, \Cref{sec:robustness} examines robustness and secondary policy comparisons, and \Cref{sec:discussion} discusses the economic interpretation and limitations. \Cref{sec:conclusion} concludes.
